% Input fragment extending finite_fibre_strengthening.tex to other fields.
\subsection{Finite-fiber compression over other fields}
\label{subsec:finite-fibre-fields}

The one-variable compression theorem does not require an ordering. The
place where rational squares were used was the conjunction of finitely
many equations. Over any field that is not algebraically closed, a
homogenized irreducible polynomial supplies the needed replacement.

\begin{lemma}\label{lem:anisotropic-binary-conjunction}
Let $K$ be a field that is not algebraically closed. There is a homogeneous
polynomial $H(A,B)\in K[A,B]$ of degree $e\ge2$ such that
\[
 H(a,b)=0\quad\Longleftrightarrow\quad a=b=0
 \qquad(a,b\in K).
\]
Consequently any finite system of polynomial equations over $K$ can be
combined into one polynomial equation without introducing variables.
For $r$ equations of degrees at most $D$, a balanced combination has
degree at most $e^{\lceil\log_2 r\rceil}D$ when $r\ge1$.
\end{lemma}
\begin{proof}
Choose a monic irreducible polynomial $h(T)\in K[T]$ of degree $e\ge2$,
and homogenize it:
\[
 H(A,B)=B^e h(A/B).
\]
This notation denotes the polynomial obtained by homogenization, also
at $B=0$. If $b\ne0$, then $H(a,b)=0$ would give a root $a/b$ of $h$
in $K$, contradicting irreducibility. If $b=0$, monicity gives
$H(a,0)=a^e$, so $a=0$. Conversely $H(0,0)=0$.

To combine residuals $f_1,\ldots,f_r$, pad the list with zero
polynomials to a power of two, and repeatedly apply $H$ to pairs.
At each stage a combined residual vanishes exactly when both its
inputs vanish. After $\lceil\log_2 r\rceil$ stages, the one remaining
equation is equivalent to the whole system. Each stage multiplies
the degree bound by at most $e$.
\end{proof}

\begin{theorem}\label{thm:finite-fibre-arbitrary-field}
Let $K$ be a field of characteristic zero that is not algebraically
closed. Let $X\subseteq\mathbb A^{d+m}_K$ be constructible over $K$,
and suppose every geometric fiber of the coordinate projection
$\pi:X\to\mathbb A^d_K$ is finite. There is a polynomial
$P(\mathbf x,U)\in K[\mathbf x,U]$ such that
\[
 \mathbf a\in\pi(X(K))
 \quad\Longleftrightarrow\quad
 \exists U\in K:\ P(\mathbf a,U)=0
 \qquad(\mathbf a\in K^d).
\]
For a number field $K$, one may take $P\in\mathcal O_K[\mathbf x,U]$,
and the construction is effective from equations and inequations for
$X$. More generally it is effective whenever the coefficient field is
presented with the algebraic operations needed for factorization,
elimination, and primitive-element computation, and an irreducible
polynomial of degree at least two is available effectively.
\end{theorem}
\begin{proof}
The algebraic part of Theorem~\ref{thm:finite-fibre-one-variable} works
over $K$. Indeed, $K$ is infinite and perfect. After passing to reduced
affine pieces and localizing on each irreducible image closure, the
generic finite algebra is a product of finite separable field extensions
of the image's function field. The primitive element theorem gives
a monic polynomial presentation. Clearing the finitely many
denominators extends the presentation over a nonempty open subset.
The complement has smaller image dimension, so induction produces
finitely many charts, including every exceptional finite fiber.

Thus, on each chart, a point lies in the image precisely when
\[
 f_1(\mathbf x)=\cdots=f_r(\mathbf x)=0,
 \qquad g(\mathbf x)\ne0,
\]
and a certain monic polynomial
\[
 h(T)=T^n+\sum_{j=1}^n\frac{N_j(\mathbf x)}{B(\mathbf x)}T^{n-j}
\]
has a root in $K$. All coefficient denominators are nonzero on the
chart; enlarge its guard to include them. For $\epsilon\in\{0,1\}$,
let $h_\epsilon(T)=h(T-\epsilon)$ and write its coefficients with
the same common denominator $B$. Define
\[
 \Psi_\epsilon(\mathbf x,U)
 =1+\sum_{j=1}^n
      N_{\epsilon,j}(\mathbf x)
      B(\mathbf x)^{2j-1}g(\mathbf x)^jU^j.
\]
When $Bg=0$ or $U=0$, this expression is $1$. Otherwise, with
$L=B^2g$, the exact identity
\[
 \Psi_\epsilon(\mathbf x,U)
   =(LU)^n h_\epsilon(1/(LU))
\]
shows that it has a zero precisely when $h_\epsilon$ has a nonzero
root. A root of $h$ gives a nonzero root of at least one of $h(T)$
and $h(T-1)$; in particular a zero root becomes the root one after
the translation. Hence these two choices cover all roots.

For each choice of chart and $\epsilon$, apply
Lemma~\ref{lem:anisotropic-binary-conjunction} to the residuals
$f_1,\ldots,f_r,\Psi_\epsilon$. This gives one equation in
$\mathbf x,U$ equivalent over $K$ to their conjunction. Take the
product of these equations over the finitely many charts and the
two translations. A finite product vanishes if and only if one
factor vanishes, which proves the assertion without adding another
witness. For an empty constructible set take $P=1$.

If $K$ is a number field, clearing a common numerical denominator
places the finitely many coefficients in $\mathcal O_K$ without
changing the zero set. All algebraic operations in the construction
are effective over a number field. The more general effective
assertion follows under the stated algorithmic hypotheses.
\end{proof}

For number fields, Lemma~\ref{lem:four-variable-anisotropic} supplies
a quadratic form in four arguments that vanishes only at the zero
tuple. A balanced tree of that form combines $r$ residuals of degree
at most $D$ with degree at most $2^{\lceil\log_4 r\rceil}D$, again
without adding variables. Thus the same witness bound is available
even for totally imaginary number fields; the failure of sums of
squares there is not an obstruction to this compression argument.

\paragraph{Why the field hypothesis cannot simply be omitted.}
Over an algebraically closed field $K$, take $X$ to be the origin in
$\mathbb A^2_K$, with the identity projection to $\mathbb A^2_K$.
All geometric fibers are finite. Its rational image nevertheless
cannot be written as
\[
 \{(a,b)\in K^2:\exists U\in K,\ P(a,b,U)=0\}
\]
for a single polynomial $P$. If $P$ has positive degree in $U$,
then every point in the nonempty open subset where its leading
$U$-coefficient is nonzero has a lift, since every nonconstant
polynomial over $K$ has a root. The image therefore cannot be a
single point. If $P$ is independent of $U$, a nonconstant polynomial
in two variables has a zero set of dimension at least one whenever
it has a zero; a constant gives either the empty set or the whole
plane. Neither possibility gives exactly the origin. Thus
non-algebraic-closedness is essential for this general one-equation,
one-additional-variable statement.
